\emph{Pour chacune des affirmations suivantes, indiquer si elle est vraie ou fausse.\\
		Chaque réponse doit être justifiée. Une réponse non justifiée ne rapporte aucun point.\\ Dans cet exercice, les questions sont indépendantes les unes des autres.}
	
	\begin{enumerate}[itemsep=5mm]
		\item On considère l'équation différentielle suivante où $y$ est une fonction de la variable réelle $x$ définie et dérivable sur $\R$.
		\[(E): \quad y'+y=2 \cos (x)\]
		
		\textbf{Affirmation 1 :} La fonction $f$ définie sur $\R$ par
		\quad $f(x)=4 \e^{-x}+\cos (x)+\sin (x)$\quad
		est solution de l'équation différentielle $(E)$.
		
		\item On considère les fonctions $f$ et $g$ définies sur $\R$ par \quad $f(x)=2 x$ \quad et \quad $g(x)=\sin (x)$.
		
		On note $\mathcal{C}_{f}$ la courbe représentative de la fonction $f$ dans un repère et $\mathcal{C}_{g}$ celle de la fonction $g$.
		
		\textbf{Affirmation 2 :} Les courbes $\mathcal{C}_{f}$ et $\mathcal{C}_{g}$ ont un seul point d'intersection.
		
		\item On considère la suite $(v_{n})$ définie pour tout entier naturel $n$ par :
		\[v_{n}=\dfrac{2 n+\sin (n)}{n+1}\]
		\textbf{Affirmation 3 :} La suite $\left(v_{n}\right)$ diverge.
		
		\item On considère la suite $(u_{n})$ définie par $u_{1}=1$ et pour tout entier naturel $n$ supérieur ou égal à 1 :
		\[u_{n+1}=u_{n}+2 n+1\]
		
		\textbf{Affirmation 4 :} Pour tout entier naturel $n$ supérieur ou égal à 1 on a $u_{n}=n^{2}$.
		
		\item On considère la suite $(u_{n})$ définie, pour tout entier naturel $n$, par $u_{n}=\e^{-n}$.
		
		On pose, pour tout entier naturel $n$,\quad $S_{n}=u_{0}+u_{1}+u_{2}+\cdots+u_{n}$\,.
		
		\textbf{Affirmation 5 :} $\lim\limits_{n \to +\infty} S_{n}=\dfrac{\e}{\e-1}$
	\end{enumerate}
